\Needspace{12\baselineskip}
\section{What this means for other arithmetic conjectures}
\label{sec:targets}

The new certified examples establish that the repository and incoming work
are useful sources of methods. They do not reduce every open problem to
matrix manipulation. The following targets have different remaining
mathematical requirements.

\subsection{A concrete nonabelian benchmark}

Let \(F=\Q(\theta)\), where \(\theta^3-\theta-1=0\). The cubic has no
rational root and has discriminant \(-23\), so its normal closure has
Galois group \(S_3\). The cubic permutation representation decomposes as
\(\mathbf1\oplus\rho\), where \(\rho\) has dimension two; the Artin
factorization is
\[
\zeta_F(s)=\zeta(s)L(s,\rho).
\]
For unramified primes, \(\operatorname{tr}\rho(\mathrm{Frob}_p)\) is the
number of roots of \(X^3-X-1\) modulo \(p\), minus one.
The identity, transposition, and three-cycle values are \(2,0,-1\).
This information is not determined by a fixed abelian congruence class.

The preceding nonabelian reports isolate one-sided prime estimates
and more elaborate nonnegative detectors for this and related fields
\cite{IncomingNonabelian,IncomingBridges}. A transparent statistic is
the smoothed Dedekind prime-ideal sum
\[
S_F(x)=\sum_{N\mathfrak p^k\le x}
\log N\mathfrak p
\left(\frac1{N\mathfrak p^k}-\frac1x\right).
\]
Its principal main term is \(\log x+c_F\).
An appropriate eventual upper estimate of
\(S_F(x)-\log x-c_F\) at square-root quality, together with the
analytic sign conditions, would give a new zero-location result.
The unresolved arithmetic task is to make the source's integral
construction record these splitting types and retain all required
uniform size, spanning, and determinant bounds.

The earlier discriminant-\(229\) analysis is useful: it identifies
an integral-frame obstruction and a possible projective-lattice remedy.
Those are local advances. Their existence alone does not imply the
needed global estimate. A sensible next milestone is one explicitly
stated uniform inequality for one fixed cubic field.

\subsection{Algebraic \texorpdfstring{\(K\)}{K}-theory and Selmer problems}

The source's cyclotomic odd-cocycle bound supplies a uniform arithmetic
integer controlling lengths. It is a promising input to finite
vanishing certificates. However,
\src{Lib/Cyclotomic/OddSha.lean}{\texttt{OddSha.lean}} explicitly separates
its cocycle model from a comparison with standard \(\Sha^1\), and does
not prove an equality of Iwasawa characteristic ideals.

To turn this into a \(K\)-group computation, one needs a correct
comparison with the relevant cohomology or Selmer group, all local
conditions and parity conventions, and the finite arithmetic certificate
in the right order of quantifiers. An abstract small length is insufficient.
The congruence diagnostic in Section~\ref{sec:congruence-smith} gives a
quick preliminary screen: rational conjugate self-duality makes the
full congruence term infinite. When that obstruction is absent,
actual full-rank relation minors can certify a finite bound.

The question is whether the arithmetic data satisfy the source theorem's
hypotheses. A matrix example with the same shape is not enough.
Accordingly, no Vandiver, general \(K_{4m}(\Z)\)-vanishing, or BSD
conclusion is asserted here.

\subsection{Why direct elliptic-curve or BSD transplantation fails}

The source \code{SelmerDatum} requires absolute residual irreducibility,
finite global and local terms, specified inertia behavior, and compatibility
of a conjugate-dual pairing with the mirror construction.
An elliptic Tate module is polarized, and its inertia and pairing data
cannot simply be substituted into an interface designed for the source's
special representations.

The previous structural report already discusses the pairing-sign
obstruction \cite{IncomingStructural}. The exact Smith-factor viewpoint
adds an operational test of finiteness, but does not remove that
obstruction. A successful application would need a theorem about
the correct signed or quotiented component and its actual local
conditions, followed by a comparison to the desired arithmetic group.

\subsection{Zero simplicity and correlations}

The endgame locates zeros by excluding poles of transforms with controlled
signs. It does not distinguish a zero of multiplicity one from a zero
of larger positive multiplicity on the allowed line. Nor do coarse
one-sided prime estimates encode pair correlation at microscopic
scales. Transferring the zero-location conclusion to simplicity,
linear independence of ordinates, or random-matrix statistics would
require new observables and estimates.

\subsection{Research priorities after this report}

\begin{center}
\begin{tabularx}{\linewidth}{@{}p{0.24\linewidth}X@{}}
\toprule
\textbf{Next problem} & \textbf{A concrete success criterion}\\
\midrule
Other cycle lengths &
Classify the allowed defect patterns and prove exact optima for an
infinite residue class outside \(8\Z\). The new \(50\)-vertex certificate
is a test case with positive holonomy.\\
Coefficient-depth growth &
Prove matching upper bounds on delayed \(s_0(n)\) using genuine
trace-kernel rows and integral elimination. The exact coefficient
threshold does not by itself control higher-term cancellation.\\
Fixed nonabelian field &
Prove a Frobenius-sensitive arithmetic inequality with constants and
uniformity sufficient for the analytic detector.\\
Finite arithmetic vanishing &
Complete the standard-Selmer comparison and supply the necessary
arithmetic determinant or Hecke certificate, with explicit coefficient rings.\\
\bottomrule
\end{tabularx}
\end{center}
